\section{Measures}
\subsection{Information}
To quantify the \textcolor{Blue}{\textbf{amount of information} in an event}
occurring with probability $p$, we find a function $\psi$ such that:
\begin{itemize}
  \item $\psi(p)$ is \textcolor{Bittersweet}{non-negative}.
  \item $\psi( 1 ) =0$. Definite events carry no information.
  \item $\psi(p)$ is \textcolor{Bittersweet}{continuous}, and decreasing in
    $p$. Unlikely events carry more information.
  \item $\psi\left( p_1p_2 \right) = \psi(p_1)+\psi(p_2)$. Independent events
    carry additive information.
\end{itemize}
The only function that satisfies these axioms is, for some base $b$:
\[\psi\left( p \right) =\log_b{\frac{1}{p}}\qquad p\in \left[ 0,1 \right].\]
\vspace{-4mm}
\begin{center}
  \includegraphics[width=0.07\textwidth]{information.png}
\end{center}
\vspace{-4mm}
We consider \textcolor{ForestGreen}{$b=2$} to measure information in
\textcolor{Bittersweet}{bits}. Any other base is equivalent up to a constant
factor.

\subsection{Shannon Entropy}
Consider a discrete random variable $X$ with PMF $p(x)$.  To quantify the
\textcolor{Blue}{\textbf{entropy}} of (the distribution of) $X$ such that:
\begin{itemize}
  \item $H(\mathbf{p})$ is \textcolor{Bittersweet}{continuous} as a function of
    $\mathbf{p}$.
  \item $H(\mathbf{p})$ is \textcolor{Bittersweet}{maximised for the uniform
    distribution}, and is increasing in the number of outcomes.
  \item Considering successive decisions, the following holds:
    \[H(\mathbf{p})=H(p_1+p_2,\ldots,p_N) + \left( p_1+p_2 \right) H_2\left(
    \frac{p_1}{p_1+p_2}\right)\]
\end{itemize}

The only function that satisfies these axioms is the \textcolor{Blue}{Shannon
entropy}, up to a constant factor:
\[H(P_X)=\mathbb{E}_{X\sim P_X}\left[ \log_2{\frac{1}{P_X(X)}} \right].\] 
Entropy is a measure of \textcolor{Bittersweet}{uncertainty/information} in
$X$.\\
\textcolor{ForestGreen}{Binary entropy}: In particular, for a
Bernoulli distribution with probability $p$, the entropy is given by the
\textcolor{ForestGreen}{binary entropy function}:
\[H_2(p)=-p\log_2{p}-(1-p)\log_2{1-p}.\]
\textcolor{ForestGreen}{Uniform distribution}: For a uniform distribution over
$N$ outcomes, the entropy is given by:
\[H(\mathbf{p})=\log_2{N}.\]

Entropy generalises to \textcolor{Blue}{joint} and \textcolor{Blue}{conditional
distributions:}
\begin{align*}
  H(X,Y) & =\mathbb{E}_{(X,Y)\sim P_{X,Y}}\left[ \log_2{\frac{1}{P_{X,Y}(X,Y)}} \right],\\
  H(X|Y) & =\mathbb{E}_{(X,Y)\sim P_{X,Y}}\left[ \log_2{\frac{1}{P_{X|Y}(X|Y)}} \right].
\end{align*}
$H(X|Y)$ can be interpreted as the \textcolor{Bittersweet}{uncertainty in $X$
after observing $Y$}.

\paragraph{Properties of Shannon entropy}
\begin{itemize}
  \item $H(X)\ge 0$. Entropy is \textcolor{Bittersweet}{non-negative}, and is
    zero iff $X$ is deterministic.
  \item $H(X)\le \log_2{|\mathcal{X}|}$, with equality iff $X$ is uniform.
  \item \textcolor{Bittersweet}{Chain rule}:
    \[H(X_1,\ldots,X_{n})=\sum_{i=1}^{n}
    H(X_i|X_1,\ldots,X_{i-1})\]
  \item $H(X|Y)\le H(X)$, with equality iff $X,Y$ are independent.\\
    \textcolor{Bittersweet}{NOTE:} $H(X|Y=y)$ could exceed $H(X)$.
  \item $H(X_1,\ldots,X_{n})\le \sum_{i=1}^{n} H(X_i)$, with equality iff
    $X_1,\ldots,X_{n}$ are independent.
  \item \textcolor{Bittersweet}{Data processing inequality}: For any
    deterministic function $f$,
    \[H(X)\ge H(f(X)).\]
  \item \textcolor{Bittersweet}{Information-preserving transformations}: For
    any bijection $f$, or if $Y$ depends on $X$ only through $f(X)$,
    \[H(X)=H(f(X)).\]
  \item \textcolor{Bittersweet}{Concavity}:
    \[H(\lambda P+(1-\lambda)Q)\ge \lambda H(P)+(1-\lambda)H(Q)\]
\end{itemize}

\subsection{KL Divergence}
To quantify the \textcolor{Blue}{\textbf{difference between two distributions}
$P$ and $Q$}, we use the \textcolor{Blue}{KL Divergence}:
\[D_{\mathrm{KL}}(P\|Q)=\mathbb{E}_{X\sim P}\left[ \log_2{\frac{P(X)}{Q(X)}} \right].\]
\begin{itemize}
  \item $D(P\|Q)$ is \textcolor{Bittersweet}{non-negative}, and is zero iff $P=Q$.
  \item \textcolor{Bittersweet}{Chain Rule}: Let $C:=D(P_{Y|X}\|Q_{Y|X}|P_X)$.
    \[D(P_{XY}\|Q_{XY})=D(P_X\|Q_X)+C,\] 
    \[C=\sum_{x} P_X(x)D\left( P_{Y|X}\left( \cdot|x \right) \|Q_{Y|X}\left( \cdot|x \right)  \right)\]
  \item \textcolor{Bittersweet}{Data processing inequality}: For any function $f$,
    \[D(P_X\|Q_X)\ge D(P_{f(X)}\|Q_{f(X)}).\]
\end{itemize}

\subsection{Mutual Information}
To quantify the \textcolor{Blue}{\textbf{``information gain''}} of observing
another variable, we use the \textcolor{Blue}{mutual information}:
\[I(X;Y)=H(X)-H(X|Y)=H(Y)-H(Y|X).\]
Mutual information generalises to \textcolor{Blue}{conditional mutual information}:
\begin{align*}
  I(X;Y|Z)&=H(X|Z)-H(X|Y,Z)\\
  &=H(Y|Z)-H(Y|X,Z)
\end{align*}
and \textcolor{Blue}{joint distributions}:
\[I(X;Y,Z)=H(X)-H(X|Y,Z).\]
Alternative forms:
\begin{align*}
  I(X;Y)&=D(P_{XY}\|P_XP_Y)\\
        &=\mathbb{E}\left[ \log_2{\frac{P_{XY}(X,Y)}{P_X(X)P_Y(Y)}} \right]\\
        &=\mathbb{E}\left[ \log_2{\frac{P_{Y|X}(Y|X)}{P_Y(Y)}} \right]
\end{align*}
\vspace{-4mm}
\begin{center}
  \includegraphics[width=0.12\textwidth]{venn-diagram.png}
\end{center}
\vspace{-4mm}
\paragraph{Properties of mutual information}
\begin{itemize}
  \item \textcolor{Bittersweet}{Non-negativity}: $I(X;Y)\ge 0$, with equality
    iff $X$ and $Y$ are independent.
  \item $I(X;Y)\le \min{\left\{ H(X), H(Y) \right\} }$ are
    \textcolor{Bittersweet}{upper bounds} on information gain.
  \item $I(X;Y)=H(X)$ iff $X$ is a \textcolor{Bittersweet}{deterministic}
    function of $Y$.
  \item \textcolor{Bittersweet}{Symmetry}: $I(X;Y)=I(Y;X)$.
  \item \textcolor{Bittersweet}{Chain rule}:
    \[I(X_1,\ldots,X_{n};Y)=\sum_{i=1}^{n} I(X_i;Y|X_1,\ldots,X_{i-1})\]
  \item \textcolor{Bittersweet}{Data processing inequality}: If $X\to Y\to Z$ 
    forms a Markov chain, then
    \[I(X;Z)\le I(X;Y).\]
  \item \textcolor{Bittersweet}{Partial sub-additivity}: If $Y_1,\ldots,Y_{n}$
    are conditionally independent given $X_1,\ldots,X_{n}$, and $Y_i$ depends
    on $X_1,\ldots,X_{n}$ only through $X_i$, then
    \[I(X_1,\ldots,X_{n};Y_1,\ldots,Y_{n})\le \sum_{i=1}^{n} I(X_i;Y_i).\]
  \item \textcolor{Bittersweet}{Convexity}: $I(X;Y)$ is convex in $P_{Y|X}$ for fixed
    $P_X$, and concave in $P_X$ for fixed $P_{Y|X}$.
\end{itemize}
